%% kapa: exam=1 round=21 year=2010 period=2 form=A subject=appraisal_law subject_ko=부동산관계법규 seq=22 number=62
%% source: Q-Net 감정평가사 기출문제 (cst00309) artlSeq=5024181 file=extracted/1008/5024181/2010년 제21회 감정평가사 1차시험 기출문제/2교시.hwp
%% conversion: kordoc 4.24.51 -> tools/gen1.py (auto); review=proofread-gpt-sol; flags=none
\begin{question}{62}
부동산 가격공시 및 감정평가에 관한 법령상 감정평가업자에 대한 과징금의 부과 및 징수에 관한 설명으로 옳은 것은?
\begin{choices}
  \item 위반행위로 인한 업무정지가 6개월 미만인 경우로서 과징금을 부과할 수 있는 때에는 과징금최고액의 100분의 20 이상 100분의 50 미만을 과징금으로 부과한다.
  \item 국토해양부장관은 감정평가법인인 감정평가업자에 대한 업무정지처분에 갈음하여 6억원 이하의 과징금을 부과할 수 있다.
  \item 과징금부과처분에 대한 이의신청은 국토해양부장관에게 하며, 그 이의신청에 대한 국토해양부장관의 결정에 이의가 있는 자는 행정심판을 제기할 수 없다.
  \item 국토해양부장관은 과징금납부의무자가 납부 기한 내에 과징금을 납부하지 아니한 경우 납부기한 경과 후 20일 이내에 서면으로 독촉하여야 한다.
  \item 국토해양부장관은 과징금이나 가산금을 납부하지 아니한 경우 비송사건절차법에 따라 징수하여야 한다.
\end{choices}
\end{question}
\answer{1}
